\documentclass{mwrep}
\usepackage{polski}
\usepackage[polish]{babel}
\usepackage[utf8]{inputenc}
\usepackage{tgbonum}
\usepackage{graphicx}
\usepackage{tabularx}
\usepackage{hyperref}

\renewcommand{\labelitemi}{$\bullet$}
\setlength{\parindent}{0cm}
\addto\captionspolish{\renewcommand{\tablename}{Tabela}}

\newcommand*{\titleGP}{\begingroup
\centering

{\large Studencka Pracownia Inżynierii Oprogramowania}\\Instytut Informatyki Uniwersytetu Wrocławskiego\par
\vspace*{16\baselineskip}

{\Large Rafał Hirsz, Karol Wieczorek, Krzysztof Wróbel\par}
\vspace*{\baselineskip}

\rule{\textwidth}{1.6pt}\vspace*{-\baselineskip}\vspace*{2pt}
\rule{\textwidth}{0.4pt}\\[\baselineskip]

{\Huge SYMULATOR LOTU}\\[0.2\baselineskip]

\rule{\textwidth}{0.4pt}\vspace*{-\baselineskip}\vspace{3.2pt}
\rule{\textwidth}{1.6pt}\\[\baselineskip]

\scshape
{\huge Analiza przypadków użycia}\par
\vspace*{2\baselineskip}

\begin{figure}[h]
\centering
\includegraphics[width=5\baselineskip]{flightsim-team-logo.pdf}
\end{figure}
\vfill

{\large Wrocław 2014}\par

\pagebreak

\endgroup}

\begin{document}

\thispagestyle{empty}
\titleGP

\begin{center}
\begin{table}[h]
\begin{center}
\caption{Historia zmian dokumentu}\label{T:Zmiany}
\vspace{3ex}
\begin{tabularx}{1\textwidth}{|l|l|l|X|}
\hline
Data & Wersja & Autor & Opis zmian \\ \hline
2014-01-10 & 1.0 & Karol Wieczorek & Utworzenie dokumentu \\
\hline
\end{tabularx}
\end{center}
\end{table}
\end{center}

\pagebreak

\tableofcontents

\chapter{Ustawienia}
\section{Wybór ustawień graficznych}
\begin{enumerate}
  \item Użytkownik klika w~przycisk ,,Ustawienia'' w~głównym menu.
  \item Użytkownik wybiera zakładkę ,,Ustawienia graficzne''.
  \item Użytkownik wprowadza jedną lub~więcej zmian:
    \begin{itemize}
      \item rozdzielczość okna,
      \item poziom szczegółów kokpitu samolotu,
      \item poziom szczegółów map,
      \item włączenie synchronizacji pionowej,
      \item ustawienia dotyczące graficznego interfejsu użytkownika.
    \end{itemize}
  \item Użytkownik naciska przycisk ,,Wprowadź zmiany''.
  \item Program wprowadza podane przez użytkownika zmiany.
\end{enumerate}

\section{Wybór ustawień dźwięku}
\begin{enumerate}
  \item Użytkownik klika w~przycisk ,,Ustawienia'' w~głównym menu.
  \item Użytkownik zakładkę ,,Ustawienia dźwięku''.
  \item Użytkownik wprowadza jedną lub~więcej zmian:
  \begin{itemize}
    \item głośność dźwięku,
    \item ilość kanałów.
  \end{itemize}
  \item Użytkownik naciska przycisk ,,Wprowadź zmiany''.
  \item Program wprowadza podane przez użytkownika zmiany.
\end{enumerate}

\chapter{Początek rozgrywki}
\section{Rozpoczęcie samouczka}
\begin{enumerate}
  \item Użytkownik naciska przycisk ,,Samouczek'' w~głównym menu.
  \item Program wyświetla listę wszystkich dostępnych scenariuszy. Scenariusze, które zostały już ukończone są oznaczone zielonym kolorem, pierwsza niewykonana lekcja białym, a~pozostałe czerwonym.
  \item Użytkownik wskazuje scenariusz, którą chciałby odbyć.
  \item Program wczytuje dane wybranego scenariusza (mapę, samolot, ustawienia, sytuację początkową).
  \item Rozpoczyna się właściwa rozgrywka.
\end{enumerate}

\section{Rozpoczęcie misji}
\begin{enumerate}
  \item Użytkownik naciska przycisk ,,Misje'' w~głównym menu.
  \item Program wyświetla listę wszystkich dostępnych misji wraz z~opisem i~szacowanym czasem ukończenia.
  \item Użytkownik wskazuje zadanie, które chciałby wykobnać.
  \item Program wyświetla ekran, na którym znajdują się wszystkie szczegóły, mapy i~diagramy potrzebne do~jej przejścia.
  \item Użytkownik zapoznaje się z~ww.~pomocami i~naciska przycisk ,,Rozpocznij misję''.
  \item Program wczytuje dane wybranej misji (mapę, samolot, ustawienia, sytuację początkową).
  \item Rozpoczyna się właściwa rozgrywka.
\end{enumerate}

\section{Rozpoczęcie swobodnego lotu}
\begin{enumerate}
  \item Użytkownik naciska przycisk ,,Tryb swobodnego lotu'' w~głównym menu.
  \item Program wyświetla ekran z~dostępnymi samolotami.
  \item Użytkownik wybiera samolot i~naciska przycisk ,,Dalej''.
  \item Program wyświetla ekran z~dostępnymi mapami oraz miejscami startu znajdującymi się na~poszczególnych mapach.
  \item Użytkownik wybiera mapę oraz~miejsce startu i~naciska przycisk ,,Dalej''.
  \item Program wyświetla ekran z~pozostałymi ustawieniami dotyczącymi:
  \begin{itemize}
    \item warunków pogodowych,
    \item ładunku samolotu,
    \item planu lotu,
    \item ewentualnych usterek samolotu.
  \end{itemize}
  \item Użytkownik dostosowuje rozgrywkę do własnych upodobań i~naciska przycisk ,,Dalej''.
  \item Program wyświetla podsumowanie konfuguracji lotu swobodnego.
  \item Użytkownik zapoznaje się z~podsumowaniem. Jeśli wszystko jest w~porządku, naciska przycisk ,,Rozpocznij grę'', w~przeciwnym wypadku naciska przycisk ,,Wstecz'' i~wprowadza poprawki.
  \item Program wczytuje odpowiednie dane (mapę, samolot, ustawienia, sytuację początkową).
  \item Rozpoczyna się właściwa rozgrywka.
\end{enumerate}

\chapter{Rozgrywka}
\section{Lot podczas samouczka //To jest słabe, ale nie umiem lepiej, liczę na protipy}
\begin{enumerate}
  \item W~momencie rozpoczęcia dowolnej z~lekcji samouczka program wyświetla dodatkowe okno na~komunikaty szkoleniowe oraz~mini-zadania do wykonania przez użytkownika.
  \item Co~pewien czas użytkownikowi pokazywany jest komunikat pozwalający zapoznać się z~pewnym elementem sztuki prowadzenia samolotu.
  \item Po~każdym komunikacie szkoleniowym użytkownik dostaje mini-zadanie do~wykonania (np. będąc jeszcze na~ziemi wychylić maksymalnie lotki i~wrócić potem do~początkowej pozycji).
  \item Po~poprawnym wykonaniu mini-zadania program zapisuje stan gry i~wyświetla kolejny komunikat.
  \item W~przypadku niepowodzenia (trudność mini-zadań wzrasta wraz z~postępami czynionymi w~samouczku) program automatycznie wczytuje ostatni zapisany stan gry.
\end{enumerate}

\section{Lot podczas wykonywania misji //To też}
\begin{enumerate}
  \item W~momencie rozpoczęcia misji program wyświetla dodatkowe okno z~celami misji (głównymi i~pobocznymi).
  \item Gdy~użytkownik wykona jeden z~celów misji, jest on~zaznaczany na~zielono.
  \item Gdy~użytkownik podejmie akcję, która uniemożliwi wykonanie któregoś z~celów, to~cel ten zostaje zaznaczony na~czerwono.
  \item Po~ukończeniu rozgrywki program pokazuje podsumowanie wyników misji i, w~przypadku pomyślnego zakończenia, nagradza użytkowika odpowiednimi trofeami.
\end{enumerate}

\section{Sterowanie za pomocą klawiatury}
\begin{enumerate}
  \item Podczas lotu użytownik naciska dowolny klawisz uwzględniony w~tabeli 2.1 z~dokumentu \cite{WYM}.
  \item Program odczytuje wciśnięty klawisz i~wykonuje odpowiednią reakcję.
\end{enumerate}

\section{Sterowanie za pomocą joysticka}
\begin{enumerate}
  \item Podczas lotu użytownik naciska dowolny ruch joysticka uwzględniony w~tabeli 2.2 z~dokumentu \cite{WYM}.
  \item Program odczytuje wciśnięty klawisz i~wykonuje odpowiednią reakcję.
\end{enumerate}

\section{Sterowanie za pomocą myszy}
\begin{enumerate}
  \item Podczas lotu użytownik klika myszą na~dowolny element w~kokpicie samolotu.
  \item Program sprawdza, który element kokpitu został użyty i~wykonuje stosowną akcję.
\end{enumerate}

\section{Kontakt z wieżą kontroli lotów}
\begin{enumerate}
  \item Użytkownik uaktywnia okno kontroli lotów poprzez kliknięcie na~nie.
  \item Program wyświetla użytkownikowi listę możliwych komunikatów.
  \item Użytkownik wybiera komunikat, który chce nadać.
  \item Komunikat zostaje wysłany do~wieży kontoli lotów.
  \item Po~chwili użytkownik dostaje odpowiedź.
  \item Jeśli odpowiedź kończy wymianę komunikatów, następuje powrót do~trybu lotu, w~przeciwnym wypadku powtarzane są~kroki~2-6.
\end{enumerate}

\section{Powiadomienia z wieży kontroli lotów //Krzysiu, sprawdź czy to ma w ogóle sens!}
\begin{enumerate}
  \item Wieża nadaje komunikat (np. o~zmianie miejsca lądowania) do~samolotu.
  \item Przy oknie kontroli lotów pojawia się ikona z~liczbą nieodebranych komunikatów.
  \item Użytkownik klika na~ikonę i~odczytuje wiadomości.
\end{enumerate}

\section{Korzystanie z biblioteki}
\begin{enumerate}
  \item W~dowolnym momencie użytkownik przechodzi do~menu i~naciska przycisk ,,Biblioteka''.
  \item Program łączy się z~internetową bazą dokumentów, a~następnie wyświetla prostą wyszukiwarkę.
  \item Użytkownik wyszukuje odpowiedni dokument w~sposób określony w~rozdziale~2.1.4 w~\cite{WYM}.
  \item Program pobiera dokument i~wyświetla go użytkownikowi.
\end{enumerate}

\begin{thebibliography}{9}
  \bibitem{WYM} R. Hirsz, K. Wieczorek, K. Wróbel: \textit{Symulator lotu: Specyfikacja wymagań}, Wrocław, SPIO IIUWr 2014.
\end{thebibliography}

\end{document}